%%% LearnPAd Document template / example using learnpad.cls class for styling
%%% 20140704, 
%%% Guglielmo De Angelis <guglielmo.deangelis@isti.cnr.it>
%%% Andrea Polini <andrea.polini@unicam.it>

\documentclass{learnpad}
\usepackage[utf8]{inputenc}
\usepackage[T1]{fontenc}
\usepackage{verbatim}
\usepackage[english,orig,printdayoff,nocleanlook]{isodate}

%%% ------------------------------------------------------
%%% ---------------- The Title
%%% ------------------------------------------------------
\title{Technology-oriented Learn PAd whitepaper}

%%% ------------------------------------------------------
%%% ---------------- The Sub-Title
%%% ------------------------------------------------------
\subtitle{First version}

%%% ------------------------------------------------------
%%% ---------------- The Name of the Deliverable
%%% ------------------------------------------------------
\deliverableno{D9.5}

%%% ------------------------------------------------------
%%% ---------------- The Authors
%%% ------------------------------------------------------
\authors{
Jean Simard (XWIKI)
}

%%% ------------------------------------------------------
%%% ---------------- The Editors
%%% ------------------------------------------------------
\editors{Jean Simard}

%%% ------------------------------------------------------
%%% ---------------- The reviewers
%%% ------------------------------------------------------
\reviewers{Andrea Polini (UNICAM)}

%%% ------------------------------------------------------
%%% ---------------- The date
%%% ------------------------------------------------------
\date{\today}

%%% ------------------------------------------------------
%%% ---------------- deliverable info
%%% ---------------- choose among : Report / Other / Prototype
%%% ------------------------------------------------------
\naturedeliverable{Report}%
%%% ------------------------------------------------------
%%% ---------------- deliverable dissemination levele
%%% ---------------- choose among the two options below:
\disseminationlevelpublic
% \disseminationlevelconfidential
%%% ------------------------------------------------------
\version{2.4}%
\contractualdeliverydate{24 January 2016}%
\actualdeliverydate{24 January 2016}%
\contributingwp{WP9}%

%%% ------------------------------------------------------
%%% ---------------- abstract
%%% ------------------------------------------------------

\abstract{In this white paper, you'll be shown a technical overview of the
\learnpad platform.  \learnpad platform is targeting learning of processes in
Public Administration with innovative approach combining model-driven learning
and collaborative learning.  This white paper will first expose the different
functionalities of the platform through a scenario that covers a variety of use
cases.  Then it will present the components and communications involved in order
to orchestrate these use cases.}

%%% ------------------------------------------------------
%%% ---------------- Keywords
%%% ------------------------------------------------------
\keywords{platform, white paper}

%%% ------------------------------------------------------
%%% ---------------- review table
%%% ------------------------------------------------------
\reviewoutline{2 Dec. 2015}{0.2}{}{}
\reviewdraft{27 Dec. 2015}{1.2}{}{}
\reviewinternal{10 Jan. 2016}{2.4}{Andrea Polini, Antonia Bertolino}{}
\reviewcandidatefinal{24 Jan. 2016}{3.0}{Antonia Bertolino}{}

\begin{document}

\frontmatter
\maketitle

%% ------------------------------------------------------
%% ---------------- document history
%% ------------------------------------------------------
\begin{history}
  \historyitem{0.1}{ToC}{Jean Simard}
  \historyitem{0.2}{ToC}{Benjamin Lanciaux, Jean Simard}
  \historyitem{1.0}{Write the Introduction}{Jean Simard}
  \historyitem{1.1}{Add Flow view and Component View}{Jean Simard}
  \historyitem{1.2}{Remove sections}{Jean Simard}
  \historyitem{2.0}{Improve after a review from XWiki}{Benjamin Lanciaux, Jean Simard}
  \historyitem{2.1}{Improve after a review from NME}{Jovaldas Januškevičius, Jean Simard}
  \historyitem{2.2}{Add graphics}{Jean Simard}
  \historyitem{2.3}{Improve with contributions}{Francesco Basciani, Antonello Calabrò, Alessio Ferrari, Francesca Lonetti, Alfonso Pierantonio, Andrea Polini, Gianni Rosa, Darius Šilingas, Jean Simard}
  \historyitem{2.4}{Improve with contributions from LINAGORA}{Tom Jorquera, Jean Simard, Sarah Zribi}
\end{history}

%%% ------------------------------------------------------
%%% ---------------- review table with the previous info
%%% ------------------------------------------------------
\reviewtable

% %%% ------------------------------------------------------
% %%% ---------------- acronyms
% %%% ------------------------------------------------------
% \begin{acronyms}
%   \acronym{CA}{Consortium Agreement}%
%   \acronym{DL}{Deliverable Leader}%
%   \acronym{DOW}{Description of Work}%
%   \acronym{EC}{European Commission}%
%   \acronym{EL}{Exploitation Leader}%
%   \acronym{GA}{Grant Agreement}%
%   \acronym{IPR}{Intellectual Property Rights}%
%   \acronym{PAB}{Project Advisory Board}%
%   \acronym{PCB}{Project Coordination Board}%
%   \acronym{PL}{Project Leader}%
%   \acronym{PMB}{Project Management Board}%
%   \acronym{PO}{Project Officer}%
%   \acronym{SL}{Scientific Leader}%
%   \acronym{S\&T}{Scientific and Technical}%
%   \acronym{TL}{Technical Leader}%
%   \acronym{WP}{Work Package}%
%   \acronym{WPL}{Work Package Leader}
% %   \acronym{\dots}{\dots~\dots}%
% \end{acronyms}

\tableofcontents

%%% ------------------------------------------------------
% In case you don't need one of the following list 
% just comment the line
%%% ------------------------------------------------------

% \listoftables 
% \listoffigures 
% \listoflistings

%%% ------------------------------------------------------

\mainmatter

%%% ------------------------------------------------------
%%% ---------------- Start with chapter and sections here!
%%% ------------------------------------------------------

\chapter{E-Learning and Knowledge Management of processes in Public Administration}
\label{ch:intro}

\section{Why \learnpad?}
In modern society public administrations (PAs) are undergoing a transformation
of their perceived role from controllers to proactive service providers, and are
under pressure to constantly improve their service quality while coping with
quickly changing context (changes in law and regulations, societal
globalization, fast technology evolution) and decreasing budgets.  Civil
servants are challenged to understand and put in action latest procedures and
rules within tight time constraints.

\section{What is \learnpad?}
\learnpad is an innovative holistic e-learning platform for PAs that enables
process-driven learning and fosters cooperation and knowledge-sharing.
\learnpad technical innovation is based on four pillars:
\begin{enumerate}
	\item a new concept of model-based e-learning (both process and knowledge)
		\begin{quote}
			Learning of processes by Civil Servants is driven by a very
			structured information: models
		\end{quote}
	\item open and collaborative e-learning content management
		\begin{quote}
			Learning should be based on open data and everybody can participate
			to improve these data
		\end{quote}
	\item automatic, learner-specific and collaborative content quality assessment
		\begin{quote}
			Every new content is verified by automatic systems and also by other
			Civil Servants, collaboratively
		\end{quote}
	\item automatic model-driven simulation-based learning and testing
		\begin{quote}
			Civil Servants can run a process in a risk-free simulation in order
			to learn and be evaluated on their new knowledge
		\end{quote}
\end{enumerate}

\section{Who will benefit \learnpad?}
\learnpad is designed with and for the Public Administration.  Civil Servants
from these Public Administrations implements processes everyday: they must be
aware of any evolution or any new process.

\section{When will \learnpad exist?}
\learnpad is a European Union funded project and started on
\printdate{01/02/2014}.  It's a two and half years long project so it will end
on \printdate{31/08/2016}.  In the meantime, a first prototype has been
evaluated by Regione Marche Administration (Italy) on \printdate{10/11/2015}.  A
second release of the platform will be release on \printdate{15/04/2016}.  A
final release is expected for \printdate{31/04/2016}.

\section{What is in this white paper?}
This white paper will present the technical solutions that have been merged
together to meet the requirements of this \learnpad platform.  In the next
section, you'll meet with the \learnpad platform.  In the last section, you'll
meet the partners that participate in the elaboration of this \learnpad
platform.

\chapter{\learnpad platform: how to use it?}
In this section, we'll provide different use cases which will give an overview
of \learnpad Platform.  These use cases are user-driven: you'll see Modeler use
cases and Learner use cases.  But first, let's explain what is the \learnpad
Platform and what are the kind of users that will use this \learnpad Platform.

\section{Overview of \learnpad Platform}
The Figure~\ref{fig:flow-view} shows the global workflow happening inside the
\learnpad platform.  The Modeling Environment component
\includegraphics[height=1.5ex,keepaspectratio]{figures/1.png} is the place where
the models are designed.  Then they are verified by the Model Verification
component \includegraphics[height=1.5ex,keepaspectratio]{figures/2.png} before
being transformed by the Transformations component
\includegraphics[height=1.5ex,keepaspectratio]{figures/3.png} in different kinds
of representations:
\begin{itemize}
	\item Wiki pages for the Collaborative Workspace component
		\includegraphics[height=1.5ex,keepaspectratio]{figures/4.png}
	\item Ontologies for the Ontology Recommender component
		\includegraphics[height=1.5ex,keepaspectratio]{figures/5.png}
	\item Business Process files for the Simulation Environment component
		\includegraphics[height=1.5ex,keepaspectratio]{figures/6.png}
\end{itemize}
Eventually, Collaborative Workspace is checking the content co-created by Civil
Servants by sending them to the Content Analysis component
\includegraphics[height=1.5ex,keepaspectratio]{figures/7.png}.

\begin{figure}[!htp]
	\centering
	\includegraphics[width=.6\paperwidth,keepaspectratio]{figures/logical-view.png}
	\caption{Flow view of the \learnpad platform}
	\label{fig:flow-view}
\end{figure}

\section{Kind of users in \learnpad Platform}
\learnpad platform has the goal to document processes in the Public
Administration for Civil Servants to learn them.  It means that the first users
of the platform will be Civil Servants: they will learn or improve their
existing knowledge about processes.

In the meantime, someone has to document these processes: the Modeler.  Based on
information of an experienced Civil Servant, the Modeler will translate the
process into models that will be injected in the \learnpad platform.

The two following sections will be about the typical use cases of a Modeler and
of a Civil Servant, and more specifically, a Learner.

\section{Modeler use cases}
\subsection{Modeling}
The first task of the Modeler is to create and document a existing process from
the Public Administration.  Working on a Modeling tool, he will model Business
Processes.  But in \learnpad, the Modeler brings a lot more than Business
Process modeling.  The Modeler will also complete Business Process model with:
Organization Structure, Documents and Knowledge, Competencies, Case Management,
KPI, Motivations.

\paragraph{Behind the scene}
The Modeler will work on the Modeling Environment component
[\ref{sec:modeling-environment}].  In \learnpad, two of these components have
been tested: Adoxx from BOC, and MagicDraw from No Magic.  Both these softwares
have implemented specific meta-models for Business Process, Organization
Structure, Documents and Knowledge, Competencies, Case Management, KPI,
Motivations.  On top of that, they also include a weaving meta-model that links
all these meta-models together.

\subsection{Publishing the models}\label{sse:PublishingModels}
Once the models has been finalized, the Modeler will publish them in order to
make them accessible to the Learners.

\paragraph{Behind the scene}
Once the Modeler decide to publish the models, the Modeling Environment
component [\ref{sec:modeling-environment}] will push the models toward the
\learnpad Core Platform component [\ref{sec:core}].  The \learnpad Core Platform
component [\ref{sec:core}] will have to check the models then import them into
the other components (more details about these processes in the next sections).

\subsection{Verification of the models}
When the Modeler is publishing the models, the \learnpad Platform may reject
them because of errors.  In fact, the \learnpad Platform will run a list of
automatic verifications to check the consistency of models.

\paragraph{Behind the scene}
In the \learnpad Platform, a Model Verification component
[\ref{sec:model-verification}] is in charge of formal verifications of the
models.  It will provide help to the Modeler by spotting errors and providing
warnings.  Once \learnpad Core Platform component [\ref{sec:core}] has check
that models don't contain errors, it will pushed them towards other component of
the \learnpad Platform, mainly Collaborative Workspace component
[\ref{sec:collaborative-workspace}], Ontology Recommender component
[\ref{sec:ontology-recommender}] and Simulation Environment component
[\ref{sec:simulation-environment}] (more details on these component in the
following sections).

\subsection{Resolve feedbacks}
Civil Servants may spot inconsistencies in the models.  In this case, they'll
fill a feedback intended for the Modeler.  When starting his modeling tool, the
Modeler will see these feedbacks, he will either fix them or reject them, then
publish the new model toward the \learnpad Platform.

\paragraph{Behind the scene}
Civil Servants can fill a feedback from the Collaborative Workspace component
[\ref{sec:collaborative-workspace}].  The Modeling Environment component
[\ref{sec:modeling-environment}] will ask feedback to the \learnpad Core
Platform component [\ref{sec:core}] which will gather them from the
Collaborative Workspace component [\ref{sec:collaborative-workspace}].  Once the
Modeler has modified and fixed the models, then the publish procedure will be
the same as already described in Section~\ref{sse:PublishingModels}.  See also
Section~\ref{sse:FormulateFeedbackToModeler} for how Civil Servants are creating
feedbacks in the Collaborative Workspace component
[\ref{sec:collaborative-workspace}].

\section{Learner use cases}
\subsection{Discover a new process}
At some point, a Learner will want to learn about a new process from the Public
Administration.  Once the Learner is logged onto the \learnpad Platform, he'll
see a list of the available processes on his dashboard.  He'll be able to
navigate through the process, going forward and backward in the series of
activities of the process and all of the contextual information linked to this
process.

\paragraph{Behind the scene}
The models are imported into the Collaborative Workspace component
[\ref{sec:collaborative-workspace}] which is wiki-based.  When the \learnpad
Core Platform is pushing models towards the Collaborative Workspace component
[\ref{sec:collaborative-workspace}], the models are transformed into wiki pages
by a Transformations component [\ref{sec:transformations}]: each element of the
model is becoming a structured wiki page.  Once the models have been imported
into the Collaborative Workspace component [\ref{sec:collaborative-workspace}],
the Learner can navigate through the models directly in the wiki.

\subsection{Getting recommendations}
In order to improve the learning experience of the Learner, the wiki will also
display recommendations to the Learner as he browses through the models in the
wiki.  Recommendations could be about people to get help from, existing example
that could help to understand the process, multimedia materials, etc.

\paragraph{Behind the scene}
In the \learnpad Platform, the Ontology Recommender component
[\ref{sec:ontology-recommender}] is able to give recommendations based on
contextual information from the browsing of the Learner.  For example,
information about who he is and what he's looking at could be used to suggest
colleagues that may help him on this particular activity of the learned process.
In order to infer recommendations from the contextual informations, Ontology
Recommender component [\ref{sec:ontology-recommender}] is transforming models
into an ontology through a Transformations component
[\ref{sec:transformations}].

\subsection{Collaborate to the documentation}
At some point, the Learner will get experience on the process and may point out
the weaknesses of the documentation.  Looking at a particular activity in the
process, he'd like to complete with some pertinent information.  He can create a
new document which will be linked to this specific activity of the process, and
write, based on his experience, new content that could help other Learners.
Note that everybody is able to collaborate so the documentation of the \learnpad
Platform is not only tight to the Modeler but also to the collaboration of every
Civil Servant in the Public Administration.

\paragraph{Behind the scene}
In the Collaborative Workspace component [\ref{sec:collaborative-workspace}],
each wiki page will display a possibility to create a collaborative content
attached to this page.  The collaborative content will create a new wiki page
that will be editable by any Civil Servant willing to improve the documentation.

\subsection{Reviewing collaborative content}
Collaborative contents give the ability to Civil Servants to cooperate in the
improvement of the documentation of the processes.  \learnpad Platform will
execute automatic text verifications in order to improve the quality of the
produces content.  The results of these executions will be shown to an
experienced Civil Servant who may improve the collaborative content accordingly
or reject the propositions.

\paragraph{Behind the scene}
\learnpad Platform provides a Content Analysis component
[\ref{sec:content-analysis}] which is able to run semantic analysis on the
textual content and detects text vagueness.  A specific Graphical User Interface
provide a review workflow inside the wiki, only available to Civil Servant
experienced on the given process.

\subsection{Formulate a feedback to Modelers}\label{sse:FormulateFeedbackToModeler}
Even if Learners can improve documentation with these collaborative content (see
previous section), it's also possible they think there is a mistake in the model
and would like to contact the Modeler in order to fix it.  Like collaborative
content, each wiki page is providing the functionality to write a feedback
directly to the Modeler.  Usually, these functionality will be restricted to any
experienced Civil Servant on the concerned process.

\paragraph{Behind the scene}
Based on the role of the user, his experience of the process, he will or not be
able to propose a feedback to Modelers.  When the Modeler will start his
modeling tool, the Modeling Environment component
[\ref{sec:modeling-environment}] will ask \learnpad Core Platform about any new
feedback to display to the Modeler.

\subsection{Simulate a process}
\learnpad Platform provides a way to run a process which is very efficient way
to learn.  The Learner will be teached by simulating a process with simulated
data; he'll be able to execute the activities of the process, do his own
mistakes.

\paragraph{Behind the scene}
From the Collaborative Workspace component [\ref{sec:collaborative-workspace}],
the Learner will be able to access a simulation of the process.  The process is
simulated by the Simulation Environment component
[\ref{sec:simulation-environment}], which provides its own Graphical User
Interface embedded into the wiki.  The Simulation Environment component
[\ref{sec:simulation-environment}] uses documents and data in order to execute a
full process.  During the simulation, two things are happening:
\begin{itemize}
	\item the context of the simulation (who, what is currently simulated, what
		are the data, etc.) is sent to the Ontology Recommender component
		[\ref{sec:ontology-recommender}] in order to provide real-time
		recommendations to the Leaner during the simulation
	\item the Learner will be monitored during the simulation (time of
		achievement, number of successes and errors, etc.) by the Monitoring
		Component which is part of the Simulation Environment component
		[\ref{sec:simulation-environment}]
\end{itemize}

\chapter{\learnpad platform: solutions and technologies}
\label{ch:platform}
In the previous section, we explained what were the functionalities of the
\learnpad Platform and how the users were interacting with it.  A lot of
different components are implied in these processes.  This section will present
the different components of the \learnpad platform.  This component view will
show which are the implemented components and how they are connected to each
other into the \learnpad platform (see Figure~\ref{fig:component-view}).

\begin{figure}[!htp]
	\centering
	\includegraphics[width=.6\paperwidth,keepaspectratio]{figures/component-view.png}
	\caption{Component view of the \learnpad platform}
	\label{fig:component-view}
\end{figure}

\section{Learn PAd Core}\label{sec:core}
\paragraph{Description}
\learnpad platform has been designed as modular as possible.  Everything is
plugged onto this \learnpad Core component.  In order to be independent as much
as possible from the technologies used in each component, each component is
communicating with the \learnpad Core component with REST APIs.

This component is developed inside the XWiki software.  This means that storage
is shared with all components that are part of XWiki.  However, it communicates
with all the components with REST which make it a replaceable component.

\paragraph{Provided services}
\begin{itemize}
	\item Orchestrator for inter-component communications
	\item Storage system for the \learnpad Platform\footnote{Note that each component may also have its own storage system.}
\end{itemize}

\paragraph{Technologies}
\begin{description}
	\item[Java] for platform server development
	\item[Jax-RS/Restlet] Framework to implement REST services
	\item[Jax-B] Framework to serialize Java class to XML
	\item[Component Manager] XWiki library for a component-oriented architecture
	\item[Jetty] Platform webserver
\end{description}

\section{Modeling Environment}\label{sec:modeling-environment}
\paragraph{Description}
Modeling Environment component is a specific component.  It's used for modeling
but once used, it can be shutdown.  However, the \learnpad Core component is
assumed to be always online.  Therefore, communication between these two
components are always a one way communication: the Modeling Environment
component send or ask information to the \learnpad Core component.

The format Modeling Environment component is pushing toward the \learnpad
platform is a ZIP archive (files \verb+*.lpzip+) which contains:
\begin{itemize}
	\item Set of models [format:XML/Adoxx or XML/MagicDraw]
	\item Images [format: BMP, JPG, PNG, SVG]
	\item Images map [format:HTML-Map]
	\item Business Process models [format:BPMN 2.0]
\end{itemize}

\paragraph{Provided services}
\begin{itemize}
	\item TO BE DONE
\end{itemize}

\paragraph{Technologies}
\begin{description}
	\item[NoTech] TO BE DONE
\end{description}

\section{Collaborative Workspace}\label{sec:collaborative-workspace}
\paragraph{Description}
The Collaborative Workspace is divided into two main components.  Since the
\learnpad Core component is not doing any transformation on the files (\learnpad
Core component is more like a router), the Collaborative Workspace component is
getting raw files.  Therefore, the Wiki component is downloading the raw files
from the \learnpad Core component, then use an EMF Transformation component (see
Section~\ref{sec:transformations}) to transform models into wiki pages.

The Wiki component is the most important part of the Collaborative Workspace
component.  The Graphical User Interface will display models.  Wiki component
will download models, ask for transformation, and import the resulting wiki
pages from the transformation.

\paragraph{Provided services}
\begin{itemize}
	\item Graphical User Interface for Civil Servants
	\item Embedding simulation in the Graphical User Interface
	\item Provide widget for recommendations to the Civil Servants
\end{itemize}

\paragraph{Technologies}
\begin{description}
	\item[Velocity/Groovy] Scripting languages for server-side rendering in the wiki
	\item[HTML5] Web technologies for interactive applications in the wiki
	\item[Bootstrap] for client-side web interface design
	\item[Java] for platform server development
\end{description}

\section{Transformations}\label{sec:transformations}
\paragraph{Description}
The models respectively, produced by the ADOxx, and Magic Draw Modelling
Environment are transformed to the corresponding Wiki pages by means of
techniques and tools provided by the Eclipse Modeling Framework (EMF). Such
models, describe the same business process, and are not directly comparable due
to the different standard representation. In order to normalize the artifacts,
UDA developed a transformational architecture which provide an Ecore
representation for them and bridge the abstraction distance between the Learn
Pad models and the wiki structures. The architecture performs several
operations, among them:
\begin{itemize}
	\item a model-to-model transformation
	\item a model-to-code transformation in order to produce the XWiki structure
		which are sent to the Wiki component.
\end{itemize}

\paragraph{Provided services}
\begin{itemize}
	\item Pre-processing of models for preparation to EMF parsing
	\item Transformation of models from one meta-model to another one
	\item Serialization of model
\end{itemize}

\paragraph{Technologies}
\begin{description}
	\item[EMF (Eclipse Modeling Framework)] is a modeling framework and code generation facility for building tools and other applications based on a structured data model
	\item[ATL (Atlas Transformation Language)] is a model-to-model transformation engine where a source model is translated into a target model Mb according to a transformation definition
	\item[Acceleo] is a model-to-code transformation engine that produce the code according to a template-based rules
\end{description}

\section{Ontology Recommender}\label{sec:ontology-recommender}
\paragraph{Description}
For the same reasons than the Collaborative Workspace component, the Ontology
Recommender component is getting raw files from \learnpad Core.  Therefore, the
Recommender component is getting these files, sending them to the XSL-T
Transformations component which return back ontologies of the models.

The Recommender component is using an ontology to organize the information of
the models.  The ontology is expressed as a RDF format.  This component will be
able to make recommendations to a Civil Servants as he browses a model or play a
simulation.  Contextual information will be used to get the best possible
recommendations, for example:
\begin{itemize}
	\item Who is the current Civil Servant?
	\item What the Civil Servant is currently looking at?
	\item In case of simulation, what are the data already entered by the Civil
		Servant?
\end{itemize}

This component is transforming models into ontologies.  The technology used is
XML specific, and very efficient at parsing XML files to transform them into
other serialization, RDF ontologies in this case.

\paragraph{Provided services}
\begin{itemize}
	\item TO BE DONE
\end{itemize}

\paragraph{Technologies}
\begin{description}
	\item[NoTech] TO BE DONE
\end{description}

\section{Simulation Environment}\label{sec:simulation-environment}
\paragraph{Description}
The simulation framework provides the subsystem where learners can simulate
Business Processes interactively and is used by one or multiple civil servant(s)
in order to learn processes. It proposes a simulation platform providing a
flexible and interoperable simulation environment with support for
collaboration. The simulation engine supports both the software emulation of the
involved parties, and the provisioning of dedicated means for gathering learners
willing to train on the business process by interacting with other learners. The
infrastructure also includes monitoring feature providing feedbacks for the
evaluation of learners, business processes, and associated learning contents.

\paragraph{Provided services}
\begin{itemize}
	\item BPMN process simulation session management, in single-user, collaborative, or mixed modes
	\item graphical interface for interacting with a simulation session
	\item Automated tasks verification and validation
	\item Simulation sessions dynamic monitoring via events notifications
\end{itemize}

\paragraph{Technologies}
\begin{description}
	\item[Java] for simulator server development
	\item[Junit] for automated testing
	\item[Activiti engine] for BPMN execution
	\item[Jetty] for simulator webserver
	\item[Resteasy] for REST API creation and consumption
	\item[Jackson] for object/json conversion between client and server
	\item[HTML5] for client-side web interface
	\item[Bootstrap] for client-side web interface design
	\item[JSON form] for client-side form generation
	\item[Glimpse] for monitoring messages processing
	\item[ActiveMQ] for monitoring messages dispatch
\end{description}

\section{Model Verification}\label{sec:model-verification}
\paragraph{Description}
The Model Verification component is in charge of verifying consistency of
models.  It runs multiple tests on models in order to identify different kind of
errors.  The Modeling Environment component is able to ask \learnpad Core
component for errors and report them to the modelers.

\paragraph{Provided services}
\begin{itemize}
	\item TO BE DONE
\end{itemize}

\paragraph{Technologies}
\begin{description}
	\item[NoTech] TO BE DONE
\end{description}

\section{Content Analysis}\label{sec:content-analysis}
\paragraph{Description}
The Content Analysis component is a Web application with a REst interface, to
analise natural language content according to six quality criteria (Correctness,
Presentation Clarity, Completeness, Simplicity, Non Ambiguity and Content
Clarity). Given a text, produced by civil servants and content managers, the
component checks for language defects and, for each defect, returns the
defective part of the text, together with a suggestion for improvement.

\paragraph{Provided services}
\begin{itemize}
	\item Detection of textual defects concerning Correctness, Presentation Clarity, Completeness, Simplicity, Non Ambiguity and Content Clarity
	\item Suggestion for defects removal
\end{itemize}

\paragraph{Technologies}
\begin{description}
	\item[LanguageTool] a class library used to check the Correctness of a sentence, in terms of grammar.
	\item[GATE] a class library implemented in Java for processing natural language. It is used to check Simplicity, Non Ambiguity and Content Clarity.
	\item[JSOUP] a Java library for working with HTML. It is used to parse HTML Content for the Presentation Clarity and Completeness.
\end{description}

\chapter{About partners}
\label{ch:partners}

\section{Consiglio Nazionale delle Ricerche}
Leading partner of the Learn PAd project, CNR is an Italian Public Research
Institute.  Learn PAd is carried out by Software Engineering and Dependable
Computing (SEDC) and Formal Methods and Tools (FMT) laboratories, that provide
the Content Analysis and Model Verification components, as well as monitoring
for simulation.

\paragraph{Competencies}
\begin{itemize}
	\item Identification of defects in natural language text
	\item Formal verification of business processes
	\item Monitoring of business processes and learning assessment
\end{itemize}

\paragraph{Contact}
\begin{description}
	\item[Stefania Gnesi] \href{mailto:stefania.gnesi@isti.cnr.it}{stefania.gnesi@isti.cnr.it} or +39 050 621 2918
\end{description}

\begin{figure}[!htp]
	\centering
	\includegraphics[width=6em,keepaspectratio]{figures/cnr.png}\par
	\url{http://www.cnr.it}
\end{figure}

\begin{figure}[!htp]
	\centering
	\includegraphics[width=6em,keepaspectratio]{figures/fmt.png}\par
	\url{http://fmt.isti.cnr.it/doku.php}
\end{figure}

\begin{figure}[!htp]
	\centering
	\includegraphics[width=6em,keepaspectratio]{figures/isti.png}\par
	\url{http://www.isti.cnr.it/}
\end{figure}

\section{BOC Asset Management GmbH}
Based in Austria, it's a leading company in Modeling.  It's providing the
software Adoxx for the Modeling Environment component.

\paragraph{Competencies}
\begin{itemize}
	\item TO BE DONE
\end{itemize}

\paragraph{Contact}
\begin{description}
	\item[TO BE DONE] \href{mailto:TO@BE.DONE}{TO@BE.DONE} or PHONE
\end{description}

\begin{figure}[!htp]
	\centering
	\includegraphics[width=6em,keepaspectratio]{figures/boc-group.png}\par
	\url{http://www.boc-group.com}
\end{figure}

\section{Linagora Grand Sud Ouest SA}
LINAGORA is a French leading company in Open Source Software and a Professional
Services Provider. Its novice support solution (OSSA, Open Source Software
Assurance) is a unique solution on a French market that provides a full line of
services available through our one-stop shop in order to secure the FLOSS
oriented information systems.

\paragraph{Competencies}
\begin{itemize}
	\item Research projects management
	\item Distributed and collaborative systems
	\item Business process related technologies
	\item Open Source projects management and dissemination
\end{itemize}

\paragraph{Contact}
\begin{description}
	\item[Jean-Pierre Lorré] \href{mailto:jplorre@linagora.com}{jplorre@linagora.com} or + 33 688 346 385
\end{description}

\begin{figure}[!htp]
	\centering
	\includegraphics[width=6em,keepaspectratio]{figures/linagora.png}\par
	\url{http://www.linagora.com}
\end{figure}

\section{No Magic Europe UAB}
Based in Lithuania, it's a leading company in Modeling.  It's providing the
software MagicDraw for the Modeling Environment component.

\paragraph{Competencies}
\begin{itemize}
	\item TO BE DONE
\end{itemize}

\paragraph{Contact}
\begin{description}
	\item[TO BE DONE] \href{mailto:TO@BE.DONE}{TO@BE.DONE} or PHONE
\end{description}

\begin{figure}[!htp]
	\centering
	\includegraphics[width=6em,keepaspectratio]{figures/nme.png}\par
	\url{http://www.nomagic.com}
\end{figure}

\section{Regione Marche}
Italian Public Administration of the Marche region, it's providing expertise on
processes in Public Administration.

\begin{figure}[!htp]
	\centering
	\includegraphics[width=6em,keepaspectratio]{figures/marche.jpg}\par
	\url{http://www.regione.marche.it}
\end{figure}

\section{Fachhochschule Nordwestschweiz}
Swiss University with expertise on ontologies.  It's working on the
recommendation engine.

\paragraph{Competencies}
\begin{itemize}
	\item TO BE DONE
\end{itemize}

\paragraph{Contact}
\begin{description}
	\item[TO BE DONE] \href{mailto:TO@BE.DONE}{TO@BE.DONE} or PHONE
\end{description}

\begin{figure}[!htp]
	\centering
	\includegraphics[width=6em,keepaspectratio]{figures/fhnw.jpg}\par
	\url{http://www.fhnw.ch}
\end{figure}

\section{Università degli Studi di Camerino}
This Italian University is having expertise in Business Process.  It's working
on formal verification of processes.

\paragraph{Competencies}
\begin{itemize}
	\item TO BE DONE
\end{itemize}

\paragraph{Contact}
\begin{description}
	\item[TO BE DONE] \href{mailto:TO@BE.DONE}{TO@BE.DONE} or PHONE
\end{description}

\begin{figure}[!htp]
	\centering
	\includegraphics[width=6em,keepaspectratio]{figures/unicam.png}\par
	\url{http://www.unicam.it}
\end{figure}

\section{Università degli Studi dell'Aquila}
The University of L’Aquila (UDA) has developed a solid research and experience
in the model-driven area (software engineering) on the metamodeling field, model
manipulation and analysis.  UDA concentrates his research in the specification
of transformations and their application to several problem domains: web
engineering, software architectures and middleware-based communication systems.

\paragraph{Competencies}
\begin{itemize}
	\item metamodeling and general model management
	\item the definition and instantiation of the Learn PAd metamodel that
		permits the specification of extended business processes in the domain
		of Public Administrations
	\item a transformational architecture where models produced in the different
		modeling environments are processed in order to generate the wiki
		structures associated to the business processes
	\item the adoption of state-of-the-art tecniques and tools for designing and
		implementing model-to-model transformations
	\item the adoption of state-of-the-art tecniques and tools for designing and
		implementing model-to-code transformations
\end{itemize}

\paragraph{Contact}
\begin{description}
	\item[Alfonso Pierantonio] \href{mailto:alfonso.pierantonio@univaq.it}{alfonso.pierantonio@univaq.it} or +39 3283332668
	\item[Francesco Basciani] \href{mailto:francesco.basciani@graduate.univaq.it}{francesco.basciani@graduate.univaq.it} or +39 3336777151
	\item[Gianni Rosa] \href{mailto:gianni.rosa@univaq.it}{gianni.rosa@univaq.it} or +39 3803288267
\end{description}

\begin{figure}[!htp]
	\centering
	\includegraphics[width=6em,keepaspectratio]{figures/univaq.png}\par
	\url{http://www.disim.univaq.it/main/index.php}
\end{figure}

\section{XWiki SAS}
French open-source company, which is developing XWiki software.  It's providing
XWiki as the base of the \learnpad platform and expertise about collaboration
and knowledge management.

\paragraph{Competencies}
\begin{itemize}
	\item Expert on XWiki software used as basis for \learnpad Platform
	\item Specialized in Collaborative Workspaces and Knowledge-Base Wikis
	\item Software development and integration
	\item Open-source software and community management
\end{itemize}

\paragraph{Contact}
\begin{description}
	\item[Jean Simard] \href{mailto:simard.jean@xwiki.com}{simard.jean@xwiki.com}
	\item[Caleb Delisle] \href{mailto:calebjamesdelisle@xwiki.com}{calebjamesdelisle@xwiki.com}
\end{description}

\begin{figure}[!htp]
	\centering
	\includegraphics[width=6em,keepaspectratio]{figures/xwiki.png}\par
	\url{http://www.xwiki.com}
\end{figure}

% ---------------------------------- Start with annexes here!
% ----------------------------------

% \annex{}

% ---------------------------------- Start EndNotes here!  
% ---------------------------------- 

% Plese use this command if and only if your text includes endnotes.
% Otherwise, comment it.

% \theendnotes

% ---------------------------------- Bibliography starts here
% ----------------------------------

% \bibliographystyle{plain}
% \bibliography{biblio}

\end{document}
